\documentclass[12pt]{article}
\usepackage[utf8]{inputenc}
\usepackage[T1]{fontenc}
\usepackage{amsmath,amssymb}
\usepackage{geometry}
\geometry{margin=2.5cm}

\begin{document}

\noindent\underline{\textbf{Applications: \ The Structure Theorem for cyclic groups}}

\bigskip
\noindent Let $(G,\cdot,e)$; $a\in G$ \quad $\langle a\rangle \overset{\text{not}}{=} \{a^h \mid h\in\mathbb{Z}\}$ = the cyclic subgroup

\noindent of $G$ generated by $a\in G$.

\[
a^h =
\begin{cases}
\underbrace{a\cdot a\cdots a}_{h \text{ times}}, & h=n\in\mathbb{N} \\
e, & h=0. \\
(a^{-1})^n, & h=-n\in\mathbb{N}
\end{cases}
\]

\noindent $\langle a\rangle \leq G$.

\bigskip
\noindent $G$ is called cyclic if $\exists\,a\in G$ such that $G=\langle a\rangle$.

\bigskip
\noindent\underline{Remark}: \ If ord $G=n<\infty \Rightarrow G$ is cyclic $\Leftrightarrow \exists\,a\in G$ such that ord$\,a=n$.

\noindent $G=\{e,a,\ldots,a^{n-1}\} = \langle a\rangle$.

\bigskip
\noindent (1) The groups $(\mathbb{Z},+,0)$, $(\mathbb{Z}_n,+,\hat 0)$, $n\geq 1$ are cyclic.

\noindent (2) For every cyclic group $G$, then $G\cong (\mathbb{Z},+,0)$ \ for $G$ infinite

\bigskip
\noindent\underline{Proof}:

\noindent $(\forall)\,h\in\mathbb{Z}$, $h=h\cdot 1 \Rightarrow \mathbb{Z} = \langle 1\rangle = \langle -1\rangle$

\noindent $\mathbb{Z}_n = \langle\hat 1\rangle$.

\bigskip
\noindent $f: \mathbb{Z}\to G=\langle a\rangle$, $a\in G$.

\noindent $f(h)=a^h$ \quad $f$ a surjective morphism

\bigskip
\noindent $f$ injective $\Rightarrow$ ker$f=\{0\}$.

\noindent $f$ not injective $\Rightarrow \{0\}\neq\text{ker}f\leq\mathbb{Z} \Rightarrow \exists!\,n\geq 1$ such that ker$f=n\mathbb{Z}$

\[
G=\text{Im}f \cong \frac{\mathbb{Z}}{\text{ker}f} =
\begin{cases}
\mathbb{Z}/\{0\} \cong \mathbb{Z}, & f \text{ injective}\\
\mathbb{Z}/n\mathbb{Z} = \mathbb{Z}_n, & f \text{ not injective}
\end{cases}
\]

\bigskip
\noindent $\mathbb{Z} \xrightarrow{\ 1_{\mathbb{Z}}\ } \mathbb{Z}$. \quad ker$\,1_\mathbb{Z} = 0_\mathbb{Z} \ \Rightarrow \ \text{Im}(1_\mathbb{Z}) = \mathbb{Z} \cong \dfrac{\mathbb{Z}}{\text{ker}\,1_\mathbb{Z}}$

\bigskip
\noindent\underline{\textbf{Cayley's Theorem}}: \ $A\neq\varnothing$, $S_A = \{\sigma \mid \sigma:A\to A \text{ bijective}\}$.

\bigskip
\noindent Let $(S_A,\circ,1_A)$ be the group of all permutations of the set $A$.

\noindent A subgroup $H\leq S_A$ is called a permutation group.

\bigskip
\noindent If $f:G\xrightarrow{\ \sim\ }G'$, \quad $x^3=e$ \ $(\forall)\,x\in G \ \Rightarrow \ x'^3=e'$ \ $(\forall)\,x'\in G'$.

\noindent $x'=f(x)$, $x\in G$, \quad $x'^3 = f(x^3) = f(e) = e'$.

\end{document}
